\documentclass[../../../include/open-logic-section]{subfiles}

\begin{document}

\olfileid[tr]{sth}{story}{blv}

\olsection{Ek: Frege'nin V. Temel Yasası}

\olref[rus]{sec} içinde Russell'ın paradoksunu, Frege'nin
\emph{Grundgesetze}'de ana çizgilerini verdiği sistem için bir sorun olarak
formüle ettiğini açıkladık. Frege'nin sistemi Naif Küme Oluşturma'nın doğrudan
bir formülasyonunu içermiyordu. Bu ekte Frege'nin sisteminin ne
\emph{içerdiğini}, bunun Naif Küme Oluşturma ve Russell Paradoksu ile nasıl
ilişkili olduğunu çok kısaca açıklayacağız.

Frege'nin sistemi \emph{ikinci derecedendi} ve bir \emph{kavramın kaplamı}
kavramını formüle etmek üzere tasarlanmıştı.\footnote{Kesin konuşursak Frege
  daha genel bir kavramı, bir fonksiyonun ``değer alanını''
  biçimselleştirmeye çalışır. Kavramların kaplamları bu daha genel kavramın
  özel bir durumudur. Ayrıntılar için bkz.
  \citet[pp.\ 8--9]{Heck2012}.} Frege'den esinlenen gösterimle
$\fregeext{x}{F(x)}$'i \emph{$F$ kavramının kaplamı} için yazacağız. Bu
aygıt bir \emph{yüklemi}, ``$F$''yi alıp (birinci dereceden) bir
\emph{terime}, ``$\fregeext{x}{F(x)}$''e dönüştürür. Frege bu aygıtla
üyeliğin şu \emph{tanımını} verdi:
\[
  a \in b =_\text{df} \exists G(b = \fregeext{x}{G(x)} \land Ga)
\]
Kabaca: $a\in b$ ancak ve ancak $a$, kaplamı $b$ olan bir kavramın kapsamına
girerse. (``$\exists G$'' niceleyicisinin ikinci dereceden olduğuna dikkat
edin.) Frege ayrıca \emph{V. Temel Yasa} olarak bilinen şu ilkeyi savundu:
$$\fregeext{x}{F(x)} = \fregeext{x}{G(x)} \liff \forall x (Fx \liff Gx)$$
Kabaca: Kavramların kaplamları ancak ve ancak aynı kaplama sahiplerse
özdeştir. (Yine hem ``$F$'' hem ``$G$'' yüklem konumundadır.) Şimdi basit bir
ilke üyeliği bir özelliği sağlamaya bağlar:

\begin{lem}[\emph{Grundgesetze}'de]\ollabel{lem:Fregeextensions}
$\forall F \forall a(a \in \fregeext{x}{F(x)} \liff Fa)$
\end{lem}

\begin{proof}
$F$ ile $a$'yı sabitleyin. O zaman $a\in\fregeext{x}{F(x)}$, üyelik tanımı
gereğince ancak ve ancak
$\exists G(\fregeext{x}{F(x)}=\fregeext{x}{G(x)}\land Ga)$; V. Temel Yasa
gereğince ancak ve ancak $\exists G(\forall x(Fx\liff Gx)\land Ga)$; temel
ikinci dereceden mantık gereğince ancak ve ancak $Fa$.
\end{proof}

Bu da Naif Küme Oluşturma'yı neredeyse doğrudan verir:

\begin{lem}[\emph{Grundgesetze}'de]
$\forall F \exists s \forall a (a \in s \liff Fa)$
\end{lem}

\begin{proof}
$F$'yi sabitleyin. \olref{lem:Fregeextensions},
$\forall a(a\in\fregeext{x}{F(x)}\liff Fa)$ sonucunu verir; buradan varoluşsal
genelleştirmeyle $\exists s\forall a(a\in s\liff Fa)$ çıkar. $F$ gelişigüzel
olduğundan sonuç geçerlidir.
\end{proof}

$F$'yi $\forall x(Fx\liff x\notin x)$ ile verilen kavram olarak alırsak
Russell Paradoksu çıkar.

\end{document}
